\section{排版与编号检查}
本页及以下各页均为排版示例。封面不显示页码，摘要从 1 开始连续编号。2026 年官方规范规定了标题、正文、行距和页码要求\cite{format2026}。三线表见表\ref{tab:layout}，公式见式\eqref{eq:identity}，图形见图\ref{fig:route}。
\subsection{数学表达}
\subsubsection{公式和符号}
对任意实数 $x$，平方展开恒等式为
\begin{equation}
  (x+1)^2=x^2+2x+1.\label{eq:identity}
\end{equation}
该恒等式用于检查上下标、括号及公式编号。它不代表赛题模型；正式论文应说明每个符号的含义、单位、定义域和公式成立条件。
\subsection{图表与交叉引用}
\begin{table}[!htbp]\centering
\caption{模板中主要排版项目}\label{tab:layout}
\begin{tabular}{lll}\toprule
对象 & 排版设置 & 示例用途\\\midrule
论文题目 & 三号黑体居中 & 检查摘要页\\
一级标题 & 四号黑体居中 & 检查章节层级\\
中文正文 & 小四宋体 & 检查字体与换行\\
页码 & 阿拉伯数字居中 & 检查连续编号\\\bottomrule
\end{tabular}\end{table}
\begin{figure}[!htbp]\centering
\begin{tikzpicture}[every node/.style={draw,minimum height=9mm,minimum width=32mm,font=\normalsize}]
\node (a) at (0,0) {明确任务};
\node (b) at (4,0) {建立数学机制};
\node (c) at (8,0) {求解与检验};
\draw[->,thick] (a) -- (b); \draw[->,thick] (b) -- (c);
\end{tikzpicture}
\caption{用于检查图形边界的流程示意}\label{fig:route}
\end{figure}
\subsection{结果解释的位置}
正文应解释表图中的主要关系及其原因，并回答原问题。对于比较、验证和应用类问题，可以沿用前文模型；不需要为了章节形式而虚构新的优化目标或模型。跨页表格示例见表\ref{tab:long}，程序排版见附录\ref{app:demo}。
\clearpage
\section{跨页表格检查}
下表行数仅用于触发真实分页，检查续表标题、重复表头和底部边界。行号是排版项目编号，不是实验数据。
\begin{longtable}{@{}p{.10\linewidth}p{.30\linewidth}p{.53\linewidth}@{}}
\caption{跨页长表排版示例}\label{tab:long}\\
\toprule 编号 & 项目 & 说明\\\midrule\endfirsthead
\multicolumn{3}{c}{表\thetable（续）}\\
\toprule 编号 & 项目 & 说明\\\midrule\endhead
\midrule\multicolumn{3}{r}{续下页}\\\endfoot
\bottomrule\endlastfoot
01 & 排版检查项目 & 本行仅检查分页、列宽及文字可读性。\\
02 & 排版检查项目 & 本行仅检查分页、列宽及文字可读性。\\
03 & 排版检查项目 & 本行仅检查分页、列宽及文字可读性。\\
04 & 排版检查项目 & 本行仅检查分页、列宽及文字可读性。\\
05 & 排版检查项目 & 本行仅检查分页、列宽及文字可读性。\\
06 & 排版检查项目 & 本行仅检查分页、列宽及文字可读性。\\
07 & 排版检查项目 & 本行仅检查分页、列宽及文字可读性。\\
08 & 排版检查项目 & 本行仅检查分页、列宽及文字可读性。\\
09 & 排版检查项目 & 本行仅检查分页、列宽及文字可读性。\\
10 & 排版检查项目 & 本行仅检查分页、列宽及文字可读性。\\
11 & 排版检查项目 & 本行仅检查分页、列宽及文字可读性。\\
12 & 排版检查项目 & 本行仅检查分页、列宽及文字可读性。\\
13 & 排版检查项目 & 本行仅检查分页、列宽及文字可读性。\\
14 & 排版检查项目 & 本行仅检查分页、列宽及文字可读性。\\
15 & 排版检查项目 & 本行仅检查分页、列宽及文字可读性。\\
16 & 排版检查项目 & 本行仅检查分页、列宽及文字可读性。\\
17 & 排版检查项目 & 本行仅检查分页、列宽及文字可读性。\\
18 & 排版检查项目 & 本行仅检查分页、列宽及文字可读性。\\
19 & 排版检查项目 & 本行仅检查分页、列宽及文字可读性。\\
20 & 排版检查项目 & 本行仅检查分页、列宽及文字可读性。\\
21 & 排版检查项目 & 本行仅检查分页、列宽及文字可读性。\\
22 & 排版检查项目 & 本行仅检查分页、列宽及文字可读性。\\
23 & 排版检查项目 & 本行仅检查分页、列宽及文字可读性。\\
24 & 排版检查项目 & 本行仅检查分页、列宽及文字可读性。\\
25 & 排版检查项目 & 本行仅检查分页、列宽及文字可读性。\\
26 & 排版检查项目 & 本行仅检查分页、列宽及文字可读性。\\
27 & 排版检查项目 & 本行仅检查分页、列宽及文字可读性。\\
28 & 排版检查项目 & 本行仅检查分页、列宽及文字可读性。\\
29 & 排版检查项目 & 本行仅检查分页、列宽及文字可读性。\\
30 & 排版检查项目 & 本行仅检查分页、列宽及文字可读性。\\
31 & 排版检查项目 & 本行仅检查分页、列宽及文字可读性。\\
32 & 排版检查项目 & 本行仅检查分页、列宽及文字可读性。\\
33 & 排版检查项目 & 本行仅检查分页、列宽及文字可读性。\\
34 & 排版检查项目 & 本行仅检查分页、列宽及文字可读性。\\
35 & 排版检查项目 & 本行仅检查分页、列宽及文字可读性。\\
36 & 排版检查项目 & 本行仅检查分页、列宽及文字可读性。\\
37 & 排版检查项目 & 本行仅检查分页、列宽及文字可读性。\\
38 & 排版检查项目 & 本行仅检查分页、列宽及文字可读性。\\
39 & 排版检查项目 & 本行仅检查分页、列宽及文字可读性。\\
40 & 排版检查项目 & 本行仅检查分页、列宽及文字可读性。\\
41 & 排版检查项目 & 本行仅检查分页、列宽及文字可读性。\\
42 & 排版检查项目 & 本行仅检查分页、列宽及文字可读性。\\
43 & 排版检查项目 & 本行仅检查分页、列宽及文字可读性。\\
44 & 排版检查项目 & 本行仅检查分页、列宽及文字可读性。\\
45 & 排版检查项目 & 本行仅检查分页、列宽及文字可读性。\\
46 & 排版检查项目 & 本行仅检查分页、列宽及文字可读性。\\
47 & 排版检查项目 & 本行仅检查分页、列宽及文字可读性。\\
48 & 排版检查项目 & 本行仅检查分页、列宽及文字可读性。\\
\end{longtable}
\section{模板适用边界}
此示例只验证排版能力。摘要篇幅、公式正确性、数据来源、结论可靠性及匿名要求，应针对实际论文重新检查。目录默认关闭；后续若官方更新格式要求，以最新正式通知为准。
\begin{thebibliography}{99}
\raggedright
\bibitem{format2026} 中国研究生数学建模竞赛组委会. 2026 年论文格式规范[EB/OL]. [2026-09-21].
\url{https://cpipc.acge.org.cn/sysFile/downFile.do?fileId=97a93e3a9e074738aea95eea986508f2}.
\end{thebibliography}
